%% LyX 2.1.4 created this file.  For more info, see http://www.lyx.org/.
%% Do not edit unless you really know what you are doing.
\documentclass[11pt, oneside]{book}
%\usepackage{mathpazo}
%\usepackage{helvet}
%\renewcommand{\ttdefault}{cmtt}
%\usepackage[T1]{fontenc}
%\usepackage[latin9]{inputenc}
\pagestyle{headings}
%\setcounter{secnumdepth}{3}
%\setcounter{tocdepth}{3}
%\setlength{\parskip}{\medskipamount}
%\setlength{\parindent}{0pt}
%\usepackage{color}
%\usepackage{babel}
%\usepackage{varioref}
%\usepackage{textcomp}
\usepackage{makeidx}
%\makeindex
%\usepackage{graphicx}
%\usepackage{nomencl}
% the following is useful when we have the old nomencl.sty package
%\providecommand{\printnomenclature}{\printglossary}
%\providecommand{\makenomenclature}{\makeglossary}
%\makenomenclature
%\usepackage[unicode=true,
% bookmarks=true,bookmarksnumbered=false,bookmarksopen=true,bookmarksopenlevel=2,
% breaklinks=true,pdfborder={0 0 0},backref=false,colorlinks=true]
% {hyperref}
%\hypersetup{pdftitle={Users Manual},
% pdfauthor={Synercon Technologies, LLC},
% linkcolor=blue,urlcolor=blue,citecolor=blue,anchorcolor=blue}
%
%\makeatletter

%%%%%%%%%%%%%%%%%%%%%%%%%%%%%% LyX specific LaTeX commands.
%\pdfpageheight\paperheight
%\pdfpagewidth\paperwidth


%%%%%%%%%%%%%%%%%%%%%%%%%%%%%% User specified LaTeX commands.
% \usepackage[pdf]{pstricks}
% \usepackage{auto-pst-pdf}
% \usepackage{pstricks}
%\usepackage{pst-all}
\usepackage{calc}
\newcommand{\degree}{$^{\circ}$}
\usepackage[utf8]{inputenc}

\usepackage{graphicx}
\usepackage{atbegshi,picture}
\usepackage{lipsum}
\usepackage{verbatim}

\usepackage{subcaption}
\usepackage{mwe}
\usepackage{hyperref}

\usepackage{textcomp}
\usepackage{fancyhdr}

%\AtBeginDocument{
%  \def\labelitemiii{\(\circ\)}
%  \def\labelitemiv{\(\diamond\)}
%}
%
%\makeatother

\newcommand*{\titleGP}{\begingroup % Create the command for including the title page in the document
	\centering % Center all text
	\vspace*{\baselineskip} % White space at the top of the page 
	\begin{figure}[tbph] 
		\centering 
		\includegraphics[width=0.15\linewidth]{../media/logos/synercon_logo_v3_only}
		\label{fig:logo}
	\end{figure}
	\vspace{4cm}
	{\Huge \textbf{Forensic Link Adapter Users Manual}}\\[2\baselineskip] % Title
	{\Huge Synercon Technologies LLC}\\[2\baselineskip]
	{\Large For Software Version 1}
	\vfill % Whitespace between editor names and publisher logo
	\textcopyright {\scshape 2015} \\[0.3\baselineskip] % Year 
	\begin{center}
		Synercon Technologies LLC
	\end{center}
\endgroup}

\setlength{\topmargin}{-.5in}
\setlength{\textheight}{9in}
\setlength{\oddsidemargin}{.125in}
\setlength{\textwidth}{6.25in}

\usepackage{titlesec}
\titleformat{\chapter}[display]
{\normalfont\bfseries\filcenter}
{\Huge\thechapter}
{1ex}
{\titlerule[2pt]
\vspace{2ex}%
\LARGE}
[\vspace{1ex}%
{\titlerule[2pt]}]

\begin{document}

%\fancyfoot[C]{\thepage}

\fancypagestyle{newstyle}{
\fancyhf{} % clear all header and footer fields
%\fancyhead[l]{\bfseries \nouppercase \rightmark} % except the center
\fancyhead[l]{}
\fancyfoot[C]{\thepage} % except the center
\renewcommand{\headrulewidth}{0pt}
\renewcommand{\footrulewidth}{0pt}}

\pagestyle{newstyle}

\thispagestyle{empty}
\titleGP
\newpage
\frontmatter
\begin{center}
\Large{\textbf{Disclamer}\\[\baselineskip]}
July, 2015
\end{center}
Disclaimer: The Synercon Technologies does not warrant the completeness
and accuracy of the information produced by its software. While best
efforts by industry experts have been made to generate information
from raw digital data, there may be untested variations, add-on packages,
errors or mistakes due to the large amount of data being processed.
It is the sole responsibility of the user of the Synercon Technologies
systems and Software to determine the fitness of the information provided
by the Software for any particular application and for each individual
case. Synercon Technologies cannot warrant any loss of functionality
due to the lack of reliable network connection to the Synercon Server.
The user is responsible for establishing a reliable Internet connection
for the Software to fully function. Synercon does not warrant that
the Equipment and/or Software will be compatible with each vehicle
model. Electronic components installed on the vehicles may be replacement
parts and not programmed to match the implementations of the Equipment,
hence Synercon cannot warrant the existence or accuracy of the data
recorded on the vehicle or the existence of a fault-free environment.
The information provided by the Software is to be used only by vehicle
manufacturer product experts, engineers, professional accident investigators
and traffic crash reconstruction specialists. The user should agree
that the use of the Software and Equipment requires an appropriate
level of skill and competence. Since Synercon cannot warrant that
the Equipment will create an environment that does not set new fault
codes in the vehicle\textquoteright s electronic control module(s)
which may alter the digital data contained therein, the User assumes
all risks in using the Equipment for the purposes of preserving a
digital record. The user agrees to verify the Equipment is working
according to this manual before using it to support an investigation
or reconstruction. Synercon Technologies, LLC assumes no liability
for incidental, consequential, special or punitive damages, including
but not limited to lost profits and loss of data. The user agrees
to indemnify and hold harmless Synercon Technologies, LLC, any of
its employees, owners, officers, and managers from any legal action
involving the use of the data produced or stored by the software.
\newpage
%\maketitle
\tableofcontents
\newpage
\mainmatter

\chapter{About this Document}
\paragraph{  }
This manual is intended to teach the operator the features and use
of the Forensic Link Adapter (FLA). This is manual is designed as
a teaching tool, rather than a reference guide. There are many acronyms
and terminology used throughout this document, the reader is adviced
to familiarize themselves with the Terminology section before reading.
There are also several accompanying documentation materials provided with
the FLA. It is suggested operators examin all documents before using
the FLA to help ensure the most data, and most accurate data is retrieved.
\section*{The Field Guide}
\paragraph{  }
The Field Guide provides an overview for use of the FLA in typical
senarios. It is intended to be used as a reference for operating the FLA in the field. It also provides a troubleshooting section for common field issues.
problems.
\section*{The ECM Quick References}
\paragraph{  }
The ECM Quick References provide a very brief step-by-step guide to
aquiring data from different ECMs. These are intended to be quick
procedures for persons familiar with the FLA and general opperations.


\chapter{Synercon Products}
\paragraph{  }
There are four main components of using Synercon Technologies products.
They will be refered to in the rest of this manual, so the reader
is advised to familiarize themselved with these components to better understand the rest of this manual.
\paragraph{  }
The 4 Synercon products are
\begin{itemize}
\item The Forensic Link Adapter (FLA)
\item The FLA Preview
\item The FLA Portal
\item The Smart Sensor Simulator (SSS)
\end{itemize}
\paragraph{  }
All of these products work together to create an easy, secure system to extract data from heavy vehicles.
or an RP1210 interface to allow the operator to use OEM software.
\begin{center}
	\includegraphics[width=.8\linewidth]{../media/fla_pictures/fla_plus_sss}
\end{center}
\section{The Forensic Link Adapter (FLA)}
\vspace{1cm}
\paragraph{  }
The FLA is the device used to extract data from ECMs. The FLA is designed to extract data from ECMs in a forensically secure way. The FLA uses the SAE J1939 and SAE J1587 standards to communicate with ECMs. The FLA can extract data from ECMs by itself for select ECMs, meaning there is no need to carry a computer with OEM software to obtain data. For ECMs that the FLA cannot directly support, the FLA can be an RP1210 interface. This allows the operator to use OEM software to collect data through the FLA.
\vspace{1cm}
\begin{center}
	\includegraphics[width=.9\linewidth]{../media/fla_pictures/angled_1}
\end{center}
\vspace{1cm}
\paragraph{  }
FLA Features
\begin{itemize}
\item Forensically secure data collection
\item Stand alone data collection for select ECMs
\item Easy to read OLED screen
\item Simple and durable two button interface
\item Built in GPS antenna
\end{itemize}
\newpage


\section{The FLA Preview}
\vspace{1cm}
\paragraph{  }
The FLA Preview provides an expanded interface to the FLA while out in the field. The FLA Preview is used to preview reports while in the field to confirm the operator has the data they wish to obtain. The FLA Preview is a website hosted by each FLA. It is accessed by
entering the IP address of the FLA into a web browser.
The FLA Preview can be accessed without an internet connection.
\vspace{1cm}
\begin{center}
\includegraphics[width=.9\linewidth]{../media/fla_preview_screenshots/main_page}
\end{center}
\vspace{1cm}
FLA Preview Features
\begin{itemize}
\item Ability to preview data in field
\item Ability to download event data from select ECMs
\item Ability to view old data that has been uploaded
\end{itemize}
\newpage

\section{The FLA Portal}
\vspace{1cm}
The FLA Portal is also referenced to as Synercon's Servers. This is
the website \url{https://fla.synercontechnologies.com}.
Data collected from the FLA is uploaded to this site, and operators
are able to view full reports here.
\vspace{1cm}
\begin{center}
\includegraphics[width=.9\linewidth]{../media/fla_portal_screenshots/user_landing_page}
\end{center}
\vspace{1cm}
FLA Portal Features
\begin{itemize}
\item Fiew full reports with all avaliable event data
\item View the FLAs last known location
\item Check subscription status
\end{itemize}
\newpage


\section{The Smart Sensor Simulator (SSS)}
\vspace{1cm}
The SSS simulates sensors of a typical vehicle for ECMs. It is used
when the FLA cannot communicate with the vehicle due to damaged wiring. The SSS reduces the chance the ECM will throw a fault compared to a standard calibration harness. SSSs are specific to a type of ECM.
\vspace{1cm}
\begin{center}
\includegraphics[width=.9\linewidth]{\string"../media/sss_pictures/SSS Connected to ADEMIII\string".jpg}
\end{center}
\vspace{1cm}
SSS Features
\begin{itemize}
\item Reduces chances of throwing new faults during download
\item Provides reliable communications to ECMs
\item Can be brought to the field to eliminate the need to remove ECMs from vehicles.
\end{itemize}


\chapter{Getting Started}
\paragraph{  }
Congratulations on purchasing the only heavy vehicle scan tool purpose built for examining vehicles in a forensic or investigative context. This section provides a guide to getting the FLA System set up for the first time.


\section{What is Included with the FLA}


\subsection{The Contents of the FLA Case}
\paragraph{  }
The FLA Case has everything operators need to collect data from heavy vehicles. The FLA case has the following items:
\begin{center}
\includegraphics[width=.9\linewidth]{../media/fla_pictures/fla_case_contents}
\end{center}


\subsection{The FLA Portal and FLA Preview}
\paragraph{  }
The subscription to the TRUCKcrypt software allows operators to access the FLA Portal and FLA Preview. These will be discussed in more detail in section XXX.


\section{Register as an Operator}
\paragraph{  }
The fla-admin should send you an email to invite you to be an operator for your FLA on the FLA Portal. The fla-admin email is initiated by the Organization administrator inviting a new operator. For more details, please see Chapter XXX. It will be neccessary to have a FLA Portal account to view any data you collect with your FLA.


\section{Opperating the FLA}
\paragraph{  }
The FLA has a 4x20 character LCD screen, along with two buttons to
allow the user to interact with the device.


\subsection{Interacting with the FLA}

There are two general types of screens the FLA will display.
\\[\baselineskip]
\noindent\begin{minipage}{0.45\textwidth}% adapt widths of minipages to your needs
\begin{center}
\textbf{Prompt Screens}
\end{center}
	The prompt screens will have one or two options on the last line, corresponding to the button on that side of the FLA.
\end{minipage}% 
\hfill% 
\begin{minipage}{0.45\textwidth} 
	\includegraphics[width=\linewidth]{../media/fla_screens/explanations/prompt_screen} 
\end{minipage}\\[\baselineskip]
\noindent\begin{minipage}{0.45\textwidth}% adapt widths of minipages to your needs
\begin{center}
\textbf{Working Screens}
\end{center}
The working screens will not have an option on the last line, but usually a percentage bar. During these screens, the FLA is preforming the operation displayed. \textbf{NOTE, during these screens the RED button can cancle the current task}.
\end{minipage}% 
\hfill% 
\begin{minipage}{0.45\textwidth} 
	\includegraphics[width=\linewidth]{../media/fla_screens/explanations/loading_screen} 
\end{minipage}\\


\subsection{Connect the FLA to the Internet}
\paragraph{  }
Connecting the FLA to the internet allows the device to get updates,
set the system time, and upload data to the FLA Portal. To connect the FLA to the internet follow these steps.
\begin{enumerate}
\item Plug the FLA into a functional ethernet port using the included ethernet cable.
\item The FLA should display an IP address on the second line of the title
screen, the screen that is displayed right after the FLA finishes
booting.
\item If the FLA displays ``Ethernet Unplugged'' or ``Finding IP Address''
for more than 60 seconds, consult the troubleshooting section.
\end{enumerate}

\subsection{Setting the Time}
\paragraph{  }
Once connected to the internet, the Forensic Link Adapter can synchronoise
with universal time using a dedicated time server. This step is only
necessary if the time displayed on the title screen is not the correct
time.
To correct the FLA time, follow these steps.
\\[\baselineskip]
\noindent\begin{minipage}{0.45\textwidth}% adapt widths of minipages to your needs
	\begin{center}
		\textbf{1}
	\end{center}
	 scroll to the System Configuration screen on the FLA, and select configure
\end{minipage}% 
\hfill% 
\begin{minipage}{0.45\textwidth} 
	\includegraphics[width=\linewidth]{../media/fla_screens/sys_conf} 
\end{minipage}\\[\baselineskip]
\noindent\begin{minipage}{0.45\textwidth}% adapt widths of minipages to your needs
	\begin{center}
		\textbf{2}
	\end{center}
	Scroll to the Update System Clock screen, and select Update
\end{minipage}% 
\hfill% 
\begin{minipage}{0.45\textwidth} 
	\includegraphics[width=\linewidth]{../media/fla_screens/sys_conf_update_clock} 
\end{minipage}\\[\baselineskip]
\noindent\begin{minipage}{0.45\textwidth}% adapt widths of minipages to your needs
	\begin{center}
		\textbf{3}
	\end{center}
	Once the FLA retrieves the current time, it will confirm the new time with the operator before setting the time. The FLA will continue to update this time, so the user does not need to imediatly confirm the time.
\end{minipage}% 
\hfill% 
\begin{minipage}{0.45\textwidth} 
	\includegraphics[width=\linewidth]{../media/fla_screens/sys_conf_update_clock_confirm} 
\end{minipage}\\


\subsection{\label{sub:Setting-the-Timezone-1}Setting the Timezone}
\paragraph{  }
The timezone for the FLA can be set without an internet connection.
To set the system timezone, follow these steps.
\\[\baselineskip]
\noindent\begin{minipage}{0.45\textwidth}% adapt widths of minipages to your needs
	\begin{center}
		\textbf{1}
	\end{center}
	 Scroll to the System Configuration screen on the FLA, and select configure
\end{minipage}% 
\hfill% 
\begin{minipage}{0.45\textwidth} 
	\includegraphics[width=\linewidth]{../media/fla_screens/sys_conf} 
\end{minipage}\\[\baselineskip]
\noindent\begin{minipage}{0.45\textwidth}% adapt widths of minipages to your needs
	\begin{center}
		\textbf{2}
	\end{center}
	Scroll to the Change FLA Timezone screen, and select "Change"
\end{minipage}% 
\hfill% 
\begin{minipage}{0.45\textwidth} 
	\includegraphics[width=\linewidth]{../media/fla_screens/sys_conf_change_tz} 
\end{minipage}\\[\baselineskip]
\noindent\begin{minipage}{0.45\textwidth}% adapt widths of minipages to your needs
	\begin{center}
		\textbf{3}
	\end{center}
	Select "Next" to view the next timezone, and select "Set" to set the timezone
\end{minipage}% 
\hfill% 
\begin{minipage}{0.45\textwidth} 
	\includegraphics[width=\linewidth]{../media/fla_screens/sys_conf_change_tz_option} 
\end{minipage}\\


\subsection{Updating the FLA}
\paragraph{  }
Updating the FLA provides the latest software, with the most features
and highest chance of getting the most data. It is good practice to
preform an update before downloading a vehicle, or before going out
in the field to download.
To preform an update, folow these steps.
\\[\baselineskip]
\noindent\begin{minipage}{0.45\textwidth}% adapt widths of minipages to your needs
	\begin{center}
		\textbf{1}
	\end{center}
	Ensure the FLA is plugged into an ethernet jack, and has an IP address.
\end{minipage}% 
\hfill% 
\begin{minipage}{0.45\textwidth} 
	\includegraphics[width=\linewidth]{../media/fla_screens/title_screen} 
\end{minipage}\\[\baselineskip]
\noindent\begin{minipage}{0.45\textwidth}% adapt widths of minipages to your needs
	\begin{center}
		\textbf{2}
	\end{center}
	Scroll to the System Configuration screen, select "Configure"
\end{minipage}% 
\hfill% 
\begin{minipage}{0.45\textwidth} 
	\includegraphics[width=\linewidth]{../media/fla_screens/sys_conf} 
\end{minipage}\\[\baselineskip]
\noindent\begin{minipage}{0.45\textwidth}% adapt widths of minipages to your needs
	\begin{center}
		\textbf{3}
	\end{center}
	The FLA will now display the Update FLA Software screen. If the FLA does not have the Update option, consult the troubleshooting section. Select "Update"
\end{minipage}% 
\hfill% 
\begin{minipage}{0.45\textwidth} 
	\includegraphics[width=\linewidth]{../media/fla_screens/update_ready} 
\end{minipage}\\[\baselineskip]
\noindent\begin{minipage}{0.45\textwidth}% adapt widths of minipages to your needs
	\begin{center}
		\textbf{4}
	\end{center}
	Confirm the update
\end{minipage}% 
\hfill% 
\begin{minipage}{0.45\textwidth} 
	\includegraphics[width=\linewidth]{../media/fla_screens/update_confirm} 
\end{minipage}\\[\baselineskip]
\noindent\begin{minipage}{0.45\textwidth}% adapt widths of minipages to your needs
	\begin{center}
		\textbf{5}
	\end{center}
	Wait for the FLA to preform the update. This can take several minutes
\end{minipage}% 
\hfill% 
\begin{minipage}{0.45\textwidth} 
	\includegraphics[width=\linewidth]{../media/fla_screens/update_diag} 
\end{minipage}\\[\baselineskip]
\noindent\begin{minipage}{0.45\textwidth}% adapt widths of minipages to your needs
	\begin{center}
		\textbf{5}
	\end{center}
	When the FLA is finished updating, it will need to shutdown and be rebooted for the updates to be applied. Select Shutdown.
\end{minipage}% 
\hfill% 
\begin{minipage}{0.45\textwidth} 
	\includegraphics[width=\linewidth]{../media/fla_screens/update_shutdown} 
\end{minipage}\\[\baselineskip]
\noindent\begin{minipage}{0.45\textwidth}% adapt widths of minipages to your needs
	\begin{center}
		\textbf{5}
	\end{center}
	After the FLA has completed the shutdown, unplug all sources of power, If the operator wishes to continue interacting with the device, wait 10 seconds before plugging in power.
\end{minipage}% 
\hfill% 
\begin{minipage}{0.45\textwidth} 
	% \includegraphics[width=\linewidth]{../media/fla_screens/update_shutdown} 
	pic of awesome unplug diagram.
\end{minipage}\\[\baselineskip]


\section{Familiarize yourself with the FLA}
\paragraph{  }
After the operator has preformed these steps, the FLA is ready for
field use. To maxamize the validity and completeness
of data collected, the operator is advised to continue reading this
manual. After the user is familiar with this, the best way to get
confortable with using the FLA is to practice. This manual, and the
other supporting documents will help guide operators through the data collection and review process.


\chapter{Using the FLA}
\paragraph{  }
The FLA is a powerful data extraction device, with many features.
This section will teach the reader about the FLA, how to use it, and the features of the FLA.


\section{Working with the Forensic Link Adapter}
\paragraph{  }
This section will explain the basics of using the FLA. This section is is intended for users with little experience operating the FLA.


\subsection{Powering Up}
\paragraph{  }
The FLA will begin to boot whenever a power is supplied to the device. This can be done in one of three ways
\begin{enumerate}
\item Using the vehicle connection cable to power the FLA using the battery
of the vehicle
\item Using the wall adapter to power the FLA using a wall outlet
\item Using the 12v cigarrette adapter to power the FLA either from a vehicle
battery, or a portable battery pack
\end{enumerate}
\paragraph{  }
The FLA can accept power from both the wall adapter or 12v cigarrett adapter and the vehicle connection cable at the same time. This is useful of an operator needs to preform several downloads in the same area.
\\[\baselineskip]
\noindent\begin{minipage}{0.45\textwidth}% adapt widths of minipages to your needs
%	\begin{center}
%		\textbf{5}
%	\end{center}
	Once the FLA is finished booting, the device will display the title screen
\end{minipage}% 
\hfill% 
\begin{minipage}{0.45\textwidth} 
	\includegraphics[width=\linewidth]{../media/fla_screens/title_screen}
\end{minipage}\\[\baselineskip]
\paragraph{  }
When this screen appears, the device is ready to opperate. If the device does not reach this screen in 60 seconds, see the troubleshooting section of this manual.
\paragraph{  }
Once an investigator has established permission to download the data from a vehicle, the key needs to be turned to the on (not start) position for the FLA to communicate with the Electronic Control Modules (ECMs).
\paragraph{  }
The FLA can be plugged into the 9-pin diagnostic connector, which is usually located below the dash on the left side or to the left of the driver's seat. If the ECM was removed, then it should be connected through a Smart Sensor Simulator that emulates a truck. 


\subsection{Shutting Down}

Faulure to shut down the FLA correctly can cause loss of data. To shutdown the FLA properly, follow these steps.
\\[\baselineskip]
\noindent\begin{minipage}{0.45\textwidth}% adapt widths of minipages to your needs
	\begin{center}
		\textbf{1}
	\end{center}
	Scroll to the Shutdown screen, and select "Shutdown"
\end{minipage}% 
\hfill% 
\begin{minipage}{0.45\textwidth} 
	\includegraphics[width=\linewidth]{../media/fla_screens/shutdown} 
\end{minipage}\\[\baselineskip]
\noindent\begin{minipage}{0.45\textwidth}% adapt widths of minipages to your needs
	\begin{center}
		\textbf{2}
	\end{center}
	Confirm the shutdown
\end{minipage}% 
\hfill% 
\begin{minipage}{0.45\textwidth} 
	\includegraphics[width=\linewidth]{../media/fla_screens/shutdown_confirm} 
\end{minipage}\\[\baselineskip]
\noindent\begin{minipage}{0.45\textwidth}% adapt widths of minipages to your needs
	\begin{center}
		\textbf{3}
	\end{center}
	Wait until the FLA indicates it is safe to remove power
\end{minipage}% 
\hfill% 
\begin{minipage}{0.45\textwidth} 
	\includegraphics[width=\linewidth]{../media/fla_screens/shutdown_wait} 
\end{minipage}\\[\baselineskip]
\noindent\begin{minipage}{0.45\textwidth}% adapt widths of minipages to your needs
	\begin{center}
		\textbf{4}
	\end{center}
	When the FLA displays this screen, it is safe to remove power
\end{minipage}% 
\hfill% 
\begin{minipage}{0.45\textwidth} 
	\includegraphics[width=\linewidth]{../media/fla_screens/shutdown_done} 
\end{minipage}\\[\baselineskip]


\section{Display Screen Menu System}
\paragraph{  }
The FLA uses the front pannel display as the primary way to interact with the unit. This section discribes the design of the menu system and how to navigate it.


\section{Menu Screen Descriptions\label{sec:Menu-Screen-Descriptions}}
\paragraph{  }
Almost all of the features of the FLA can be accessed through the
main menu. The only two functions that are not avaliable are the FLA Support Mode, and previewing data, which are soly avaliable on the FLA Preview. Once the FLA boots, the device will enter the main menu, which contains the most used functions, such as running a scan, and uploading data. The below figure describes the layout of the FLA menu system.
\begin{center}
\includegraphics[width=\linewidth]{../media/graphics/menu_system}
\end{center}
\paragraph{  }
The menu system of the FLA is dynamic, it will only give the operator the option to preform an action if the FLA detects the action is possible. i.e. the FLA will not allow the operator to initiate a vehicle scan if the FLA does not detect any network traffic.


\subsection{Main Menu}
\vspace{\baselineskip}
\noindent\begin{minipage}{0.45\textwidth}% adapt widths of minipages to your needs
	\begin{center}
		\textbf{Title Screen}
	\end{center}
	This is the first screen the FLA will display when it finishes booting. This screen provides the user with an overview of the FLA's state.
\end{minipage}% 
\hfill% 
\begin{minipage}{0.45\textwidth} 
	\includegraphics[width=\linewidth]{../media/fla_screens/explanations/title_screen} 
\end{minipage}\\[\baselineskip]
\noindent\begin{minipage}{0.45\textwidth}% adapt widths of minipages to your needs
	\begin{center}
		\textbf{Shutdown Screen}
	\end{center}
	This screen will allow the user to safely shutdown the FLA.
\end{minipage}% 
\hfill% 
\begin{minipage}{0.45\textwidth} 
	\includegraphics[width=\linewidth]{../media/fla_screens/shutdown} 
\end{minipage}\\[\baselineskip]
\noindent\begin{minipage}{0.45\textwidth}% adapt widths of minipages to your needs
	\begin{center}
		\textbf{Upload Data Screen}
	\end{center}
	This screen allows the operator to upload data on the FLA to the FLA Portal.
	If the Upload option is not avaliable, see the troubleshooting section.
\end{minipage}% 
\hfill% 
\begin{minipage}{0.45\textwidth} 
	\includegraphics[width=\linewidth]{../media/fla_screens/explanations/upload_ready} 
\end{minipage}\\[\baselineskip]
\noindent\begin{minipage}{0.45\textwidth}% adapt widths of minipages to your needs
	\begin{center}
		\textbf{View Faults Screen}
	\end{center}
	This screen allows the operator to view active faults on the J1939 network
\end{minipage}% 
\hfill% 
\begin{minipage}{0.45\textwidth} 
	\includegraphics[width=\linewidth]{../media/fla_screens/explanations/view_faults} 
\end{minipage}\\[\baselineskip]
\noindent\begin{minipage}{0.45\textwidth}% adapt widths of minipages to your needs
	\begin{center}
		\textbf{GPS Screen}
	\end{center}
	This screen displays the last GPS coordinates, and the FLA system time they were recieved.
\end{minipage}% 
\hfill% 
\begin{minipage}{0.45\textwidth} 
	\includegraphics[width=\linewidth]{../media/fla_screens/explanations/gps} 
\end{minipage}\\[\baselineskip]
\noindent\begin{minipage}{0.45\textwidth}% adapt widths of minipages to your needs
	\begin{center}
		\textbf{System Configuration}
	\end{center}
	This screen allows the user to enter the system configuration menu of the FLA.
\end{minipage}% 
\hfill% 
\begin{minipage}{0.45\textwidth} 
	\includegraphics[width=\linewidth]{../media/fla_screens/explanations/sys_conf} 
\end{minipage}\\[\baselineskip]


\subsection{System Configuration Menu}
\vspace{\baselineskip}
\noindent\begin{minipage}{0.45\textwidth}% adapt widths of minipages to your needs
	\begin{center}
		\textbf{Update Software Screen}
	\end{center}
	This screen allows the operator to update the FLA software to the latest version.
	If the update option is not avaliable, see the troubleshooting section.
\end{minipage}% 
\hfill% 
\begin{minipage}{0.45\textwidth} 
	\includegraphics[width=\linewidth]{../media/fla_screens/explanations/update_ready} 
\end{minipage}\\[\baselineskip]
\noindent\begin{minipage}{0.45\textwidth}% adapt widths of minipages to your needs
	\begin{center}
		\textbf{DHCP Service Configuration Screen}
	\end{center}
	This screen allows the operator to enable or disable the DHCP services for the FLA.
\end{minipage}% 
\hfill% 
\begin{minipage}{0.45\textwidth} 
	\includegraphics[width=\linewidth]{../media/fla_screens/explanations/sys_conf_dhcp_diss} 
\end{minipage}\\[\baselineskip]
\noindent\begin{minipage}{0.45\textwidth}% adapt widths of minipages to your needs
	\begin{center}
		\textbf{Update FLA System Time Screen}
	\end{center}
	This screen allows the operator to update the FLA System time.
	If the Update option is not avaliable, see the troubleshooting section.
\end{minipage}% 
\hfill% 
\begin{minipage}{0.45\textwidth} 
	\includegraphics[width=\linewidth]{../media/fla_screens/explanations/upload_ready} 
\end{minipage}\\[\baselineskip]
\noindent\begin{minipage}{0.45\textwidth}% adapt widths of minipages to your needs
	\begin{center}
		\textbf{Change FLA Timezone Screen}
	\end{center}
	This screen allows the operator to change the FLA system timezone.
\end{minipage}% 
\hfill% 
\begin{minipage}{0.45\textwidth} 
	\includegraphics[width=\linewidth]{../media/fla_screens/explanations/sys_conf_change_tz} 
\end{minipage}\\[\baselineskip]
\noindent\begin{minipage}{0.45\textwidth}% adapt widths of minipages to your needs
	\begin{center}
		\textbf{MAC Address Screen}
	\end{center}
	This screen displays the MAC Address of the FLA, as well as the IP Address.
\end{minipage}% 
\hfill% 
\begin{minipage}{0.45\textwidth} 
	\includegraphics[width=\linewidth]{../media/fla_screens/explanations/sys_conf_mac} 
\end{minipage}\\[\baselineskip]


\subsection{Menu Screen Dynamic Lines}

Certian lines of menu screens are dynamic, and display the status
of various components of the FLA.


\subsubsection{IP Address Line}

The IP address Line can have the following states:
\begin{description}
\item [{IP: XXX.XXX.XXX.XXX}] A valid IP address. XXX can be any number
between 0 and 255. An example would be 192.168.2.45
\item [{IP: 10.0.0.1}] The IP address when the FLA's DHCP server is activated.
See the DHCP section for more information.
\item [{Ethernet Unplugged}] The FLA cannot detect a live internet connection.
If this screen persists for more than 60 seconds on a known good ethernet
port, refer to the troubleshooting section.
\item [{Finding IP Address}] The FLA detected a live internet connection,
and is searching for an IP address. If this screen percists for more
than 60 seconds on a known good ethernet port, refer to the toubleshooting
section.
\end{description}

\subsubsection{Vehicle Connection Line}

The vehicle connectoin line displays the status of the vehicle busses.
This line can have the folowing states:
\begin{description}
\item [{No Vehicle Net / Check Ign. Key}] The FLA is not able to detect
any traffic on the vehicle networks. If this screen percists for more
than 10 seconds, refer to the troublshooting section.
\item [{CAN Bus only}] The FLA only detects traffic on the CAN bus
\item [{J1708/1587 only}] The FLA only detects traffic on the J1708/1587
bus
\item [{CAN and J1708}] The FLA detects traffic on both the CAN and J1708
busses
\end{description}

\chapter{The FLA Preview}
\paragraph{  }
The FLA Preview is a website that each FLA hosts. This allows operators
to interact with the FLA in ways the front pannel are not conducive
for, such as previewing reports.


\section{Accessing the FLA Preview}
\paragraph{  }
Accessing the FLA Preview is done by entering the IP address of the
FLA into a web browser. Please see the requirements section for more
information about supported web browsers. As shown below, enter the
IP Address exactly as it appears on the IP Address line.
\\[\baselineskip]
\noindent\begin{minipage}{0.45\textwidth}% adapt widths of minipages to your needs 
	%\begin{center} 
	%	\textbf{Correct:}\\[\baselineskip] 
	%\end{center} 
	\includegraphics[width=\linewidth]{../media/fla_preview_screenshots/url_correct} 
\end{minipage}% 
\hfill% 
\begin{minipage}{0.45\textwidth} 
	\includegraphics[width=\linewidth]{../media/fla_screens/title_screen_no_veh_net} 
\end{minipage}
\\[\baselineskip]
This will load the FLA Preview main page, shown below.
\begin{center}
\includegraphics[width=.9\linewidth]{../media/fla_preview_screenshots/main_page}
\end{center}
\paragraph{  }
This page provides an overview of the current status of the device. The At a glance section will display the number of pending data packages, that is the number of data packages that have been downloaded to the FLA but not yet uploaded to the FLA Portal, and if there is a download
in progress.


\section{Using the FLA Preview}
\paragraph{  }
The FLA Preview has many uses, including allowing operators to preview data in the field, download the User Manuals, download RP1210 drivers, and many other useful tasks we will cover in this section.


\subsection{View Local Data}
\paragraph{  }
Under the View Local Data page, the operator is able to view data
on the FLA. There are two sections to local data, Pending Uploads,
and Archived Data. Pending Uploads will show any data packages that
are ready to be uploaded to the FLA Portal.
\begin{center}
\includegraphics[width=.9\linewidth]{../media/fla_preview_screenshots/local_data}
\end{center}

This page will give an overview of each pending data package. The
OEM data column will be YES if the FLA was able to recover some OEM
data, and NO if the FLA could not. The yellow Archive button will
remove the data package from being uploaded, and move it to the archived section. If an archive is unintentionally archived, the FLA will have to be serviced to retrieve the data. For more information about servicing the FLA, see the Service section.
\paragraph{  }
When a download is in progress, the row for that entry will indicate the download is in progress, and will not provide the ``View Data'' button. Once the user finishes the download, the operator can refresh the page, and will see the ``View Data'' button. For more information about when a download is finished, and how to preform a download, refer to the Field Guide.
\paragraph{  }
Selecting the ``View Data'' will provide the user with a preview of the report. This preview includes all of the standards-based information, the components found, and any identification information, as well as any in-field non-standards-based previews. Non-standards-based information is any data that is provided by OEM software, such as DDEC reports' XTR files, and CAT Snapshots. The operator will get a pireview of this data here in accordance with the type of ECM being analyzed. The full reports will be avaliable through the FLA Portal when the operator uploads the data. Typically, the following will be avaliable:
\\[\baselineskip]
\noindent\begin{minipage}{0.45\textwidth}% adapt widths of minipages to your needs 
	\begin{center} 
		\textbf{Preview of DDEC ECMs}\\[\baselineskip] 
	\end{center} 
	If non-standards-based data is collected from a DDEC ECM, an XTR file will be avaliable for download.
\end{minipage}% 
\hfill% 
\begin{minipage}{0.45\textwidth} 
	\includegraphics[width=\linewidth]{../media/fla_preview_screenshots/xtr_download}
\end{minipage}\\[\baselineskip]
\noindent\begin{minipage}{0.45\textwidth}% adapt widths of minipages to your needs 
	\begin{center} 
		\textbf{Preview of CAT ECMs}\\[\baselineskip] 
	\end{center}
	If non-standars-based data is collected from a CAT ECM, a CSV file will be avaliable containing 5 frames around the trigger for any quick stop records.
\end{minipage}% 
\hfill% 
\begin{minipage}{0.45\textwidth} 
	\includegraphics[width=\linewidth]{../media/fla_preview_screenshots/csv_download}
\end{minipage}
\\[\baselineskip]
\paragraph{  }
The preview will also contain a SHA 256 Sums file for the operator
to download. This file contains the SHA 256 sums of all data files
the FLA has created. The operator can download this file, and confirm
the data uploaded to the FLA Portal was the same data the FLA created.
For more information on SHA 256 sums, as well as recommended use,
see the SHA 256 sums of suggested practices.
\paragraph{  }
The operator can also click on the ``Archived Data'' button on the
View Local Data page to view data packages that were archived. Data
packages that are uploaded to the FLA Portal are automatically archived.
From here the user can access all of the same downloads and information
as when the data package was pending uploading. Meaning, once the
data is uploaded, the user can still downloa XTR, CSV, or SHA 256
Sum files, as well as any standards-based data aquired. If the FLA
has a large number of archived data packages, the archived page may
take several moments to load.
\begin{center}
\includegraphics[width=.9\linewidth]{../media/fla_preview_screenshots/archived_data}
\end{center}


\chapter{Using the FLA Portal}
\paragraph{  }
The FLA Portal is where operators can view full reports of the data
they collected in the field. Administrators will also use the FLA
Portal to manage FLAs, and users.


\section{Accessing the FLA Portal}
\paragraph{  }
The FLA Portal is accessed by navigating to \href{https://fla.synercontechnologies.com/}{https://fla.synercontechnologies.com/}
in a supported web browser. For a list of supported web browsers, refer to the
XXX section.
\paragraph{  }
The FLA Portal will require the operator to login, using the username and password used when the operator created the account. Note that BOTH username and password are case-sensative.
\begin{center}
\includegraphics[width=0.8\textwidth]{../media/fla_portal_screenshots/login_screen}
\end{center}

Once logged in, the operator will be able to view data aquired from
the assigned FLA.


\section{Using the FLA Portal}
\paragraph{  }
For operators, the main use of the FLA Portal is to view full reports. Administrators should refer to the Administrative Guide for further information about administrating on the FLA Portal.
\paragraph{  }
Once an operator uploaded a data package to the FLA Portal, the data will be avaliable for viewing by everyone in the group. A Logged-in operator will be able to view all data packages uploaded to the server.

\begin{center}
\includegraphics[width=0.9\textwidth]{../media/fla_portal_screenshots/user_landing_page}
\end{center}

Selecting one of these reports will open the full report.

\begin{center}
\includegraphics[width=0.9\textwidth]{../media/fla_portal_screenshots/full_report/top}
\end{center}

The FLA full report has many sections, including those listed in the next sections.


\subsection{Vehicle Information}

This section contains any components that responded to a component
Identification request. The VIN, Vehicle Odometer, and Engine Hours
will also be displayed under Engine \#1 for the appropriate network.
This section also contains the time records, including vehicle time
if avaliable, and time of the last GPS Lock.

\begin{center}
\includegraphics[width=0.9\textwidth]{../media/fla_portal_screenshots/full_report/vehicle_info}
\end{center}


\subsection{Data sections}
\paragraph{  }
Data sections, such as Historical Data, Event Data, etc., are sections that contain either standards-based data, or non-standards-based data. These sections contain all of the data the FLA has decoded. Due to the wide variety of ECMs, and ECM configurations, these sections may contain different data depending on the ECM being analyzed.


\subsection{Nomenclature}
\paragraph{  }
This section contains a breif overview of the nomenclature used in
the report.


\subsection{Network Logs}
\paragraph{  }
This section contains the logs of all of the FLA's interactions with
the ECMs. These logs are provided as a means to audit the messages
being sent. If the RP1210 passthrough is engaged, separate logs for
any passthrough transactions will be avaliable.

\begin{center}
\includegraphics[width=0.9\textwidth]{../media/fla_portal_screenshots/full_report/network_logs}
\end{center}



\subsection{Version Information}

This section provides the software versions that were used when the
report was generated.

\begin{center}
\includegraphics[width=0.9\textwidth]{../media/fla_portal_screenshots/full_report/version_info}
\end{center}


\subsection{SHA 256 Sums}
\paragraph{  }
This section provides the SHA 256 Sums for all of the data files the
FLA generated when the download was preformed. These sums are encryped
and sent to the FLA Portal for verification. The FLA Portal will generate
its own SHA 256 Sums of the data files uploaded by the FLA. The FLA
Portal will then compare the two, and verify the integrety of the
upload.

\begin{center}
\includegraphics[width=0.9\textwidth]{../media/fla_portal_screenshots/full_report/valid_shas}
\end{center}
If there is a SHA 256 Mismatch, this may suggest the data has been
tampered with.

\begin{center}
\includegraphics[width=0.9\textwidth]{../media/fla_portal_screenshots/full_report/invalid_shas}
\end{center}


\subsection{Access Information}

This section provides an audit of who is currently accessing the report,
and the device's IP address.


\chapter{Administrator's Guide}

This section is intended to teach administrators how to manage FLAs
on the FLA Portal.


\section{FLA Portal terminology}

The FLA Portal has the following categories:
\begin{description}
\item [{Organizations}] are the top level with a site administrator as
the person in charge. Typically this is a highway patrol, a state
police, or a company. The person in charge of the organization has
the ability to oversee the data and use of the different FLAs.
\item [{Groups}] are a sub-level within an organization. For example, East
District and West District may be groups within an State Police organization.
If an organization only has a few operators and a few FLAs, it may
not make sense to have different groups. 
\item [{Operators}] are trained users of the FLA and are responsible for
collecting data in the field. Operators can be assigned to different
groups. Every operator in the group is able to see the data packages
from all the FLAs assigned to the group, regardless of which operators
downloaded the data.
\item [{Forensic~Link~Adapters}] are the hardware devices that an organization
owns.
\end{description}

\section{Account Overview}

Forensic Link Adapters are assigned to each organization at the time
of purchase. Once the organization completes the purchase, the serial
numbers of each FLA will be loaded into the Organization's profile.
An invitation to set up the Organization will be sent to the site
administrator. For example, the Lieutenant in charge of the crash
reconstruction division will become the site administrator once he
or she registers by following the link sent in the invitation e-mail.

\begin{figure}
\includegraphics[width=1\textwidth]{../media/fla_portal_screenshots/FLAPortalAccountOverviewScreenShot}

\caption{\label{fig:Example-FLA-Account-overview}Example Forensic Link Adapter
Portal web page for the Account overview}
\end{figure}


Once logged in as a Site Administrator, an account overview page will
be available, much like the one shown in Figure \ref{fig:Example-FLA-Account-overview}.
Many of the fields on the web page are linked to various functions.
The top bar (in black) has links to your default login page, the latest
FLA download if you are an operator, a list of all the organizations,
FLAs, and some account access settings. The web page is set up in
three distinct tables as seen in Figure \ref{fig:Example-FLA-Account-overview}.
\begin{enumerate}
\item Organization Groups
\item FLA Operators
\item Forensic Link Adapter (FLA)
\end{enumerate}

\section{Group Management}

It is the responsibility of the Organization administrator to set
up groups and operators according to the desired structure of the
Organization. You can have operators that have access to more than
one group.

\begin{figure}
\includegraphics[width=1\textwidth]{../media/fla_portal_screenshots/FLAPortalFieldOperatorsScreenShot}

\caption{\label{fig:Example-Group-Details}Example Group Details web page from
the Forensic Link Adapter Portal.}
\end{figure}



\subsection{Adding an Operator}

There are two ways to get an operator access to the group's FLAs.
The first is to invite a new operator and the second is to invite
an existing operator. However, the existing operator must have accepted
an invitation from the Organization under a unique e-mail address.
By clicking on the lower left box the says \textsf{Invite new operator}
shown in \ref{fig:Example-Group-Details}, a web page similar to
the one shown in \ref{fig:New-Operator-invitation}. The required
user information includes name and email address. The FLA portal uses
e-mail addresses to identify operators. Therefore, if a person belongs
to different organizations, they would have to use a different e-mail
for each organization. New operators will get an e-mail from \texttt{fla-admin@synercontechnologies.com}
with a link to register with the site. Many times this e-mail will
be filtered as junk, so advising the recipient to check their e-mail
filters and junk folder may be necessary.

The Work Address details on the Invite and operator page are used
to fill in a report with 

\begin{figure}
\includegraphics[height=10cm]{../media/fla_portal_screenshots/FLAPortalInviteNewOperatorScreenShot}

\caption{\label{fig:New-Operator-invitation}New Operator invitation form on
the FLA portal website.}
\end{figure}



\subsection{Assigning an FLA to the Group}

Each group needs to have at least one Forensic Link Adapter assigned
to it. To assign an FLA to a group, click the button to see a list
of available FLAs for your organization. Only people in the group
can see the data on the FLA assigned to that group

\begin{figure}
\includegraphics[width=1\textwidth]{../media/fla_portal_screenshots/FLAPortalAssignedFLASuccessScreenShot}

\caption{Assigned FLAs}
\end{figure}



\chapter{Obtaining Data}
\paragraph{  }
The FLA can obtain standards-based data from all ECMs that support the J1587 and/or J1939 standards. The FLA can also obtain additional data, such as event data, from specific ECMs.

\section{Supported ECMs}
\paragraph{  }
The FLA can obtain additional information from the following ECMs
\begin{itemize}
\item CAT
\begin{itemize}
\item ADEM 3
\item ADEM 4
\end{itemize}
\item DDEC
\begin{itemize}
\item DDEC IV
\item DDEC V
\item Some DDEC VI
\item Some XXX/XXX
\end{itemize}
\end{itemize}
It is important to use the FLA to obtain data whenever possible. The OEM software is typically designed with maintanance in mind, not forensic soundness. Most OEM software will allow the user to erase data, while the FLA prevents accidental data loss.

\section{Using Manufacturer's Software}
\paragraph{  }
The FLA is still able to obtain some data from most non-supported ECMs, as long as they support the J1587 and J1939 standards. Additional information is typically avaliable through the manufacturers (OEM's) software. The FLA can opperate as a passthrough device, allowing an operator to use a computer with the specific ECM's software to gather additional data.
\paragraph{  }
OEM software varies between manufacturer, so different steps must be taken to collect data depending on the ECM type. Synercon provides Quick References for the most popular ECMs. These quick references will provide guidance for using OEM software.

\chapter{Suggested Practices}
\paragraph{  }
This section is intended to provide a very brief overview of using
the FLA, for a more indebth guide, refer to the Field Guide.


\section{Connecting the FLA to the Internet}

In order to upload data packages to the FLA Porta, recieve FLA updates,
or update the FLA system time, the FLA needs to be connected to the
internet.
\begin{enumerate}
\item Plug the ethernet cable from the FLA into a known good wired Internet
connection that is on the same network as your computer.
\item Navigate to the FLA Title Screen \ref{enu:Upload-to-Server}. Line
2 should begin with ``IP:'' followed by 4 sets of numbers separated
by a ``.''
\item Navigate to the FLA Upload screen. If the green option is ``Upload'',
then the FLA has successfully connected to the internet. If not, refer
to the troubleshooting section.
\end{enumerate}

\section{Connect to the FLA Preview}

There are two different ways to connect to the FLA Preview, the in-office
way, and the in-field way. The in-office method requires the internet
connection outlined in the previous section, and is only required
if the operater needs the FLA Internet connection, i.e. to open a
support channel.


\subsection{In-Office Connection}
\begin{enumerate}
\item Connect the FLA to the internet as in the above section.
\item Using a web browser on a computer, navigate to the IP Address displayed
on the FLA. If you are unable to view the page, refer to the troubleshooting
section.
\end{enumerate}

\subsection{In-Field Connection}
\begin{enumerate}
\item If the computer being used is connected to a WiFi network, disconnect
from the network before proceeding.
\item Plug the Ethernet cable from the FLA to your computer's Ethernet port.
A USB to Ethernet adapter is available in the FLA kit if your computer
does not have an ethernet jack.
\item Navigate the FLA menu to the DHCP enable screen (Screen \ref{enu:DHCP})
and enable the DHCP Server on the FLA.

\begin{enumerate}
\item The IP Address of the FLA should now be 10.0.0.1
\end{enumerate}
\item Using a web browser on a computer, navigate to the IP Address displayed
on the FLA (10.0.0.1). If you are unable to view the page, refer to
the troubleshooting section.
\item When finished, do not forget to disable the DHCP service on the FLA.
Leaving the DHCP service on could cause problems if the FLA is plugged
into an office or home network.
\end{enumerate}

\section{Initiating FLA Support mode}

In order to provide the quickest and most complete support, it may
be necessary to place the FLA in support mode. This allows Synercon
Support staff to create a secure connection to an FLA anywhere with
an internet connection. It may also be refered to as a ``Support
Channel''.

To place the FLA in support mode, or ``Open a Support Channel'',
follow these steps
\begin{enumerate}
\item Verify a Synercon Support representative is ready for you to initiate
Support Mode.
\item Connect the FLA to the internet as described in section XXX.
\item Navigate to the FLA Preview site as described in section XXX.
\item Select the ``Support'' option from the menu at the top.
\item Click the ``Open the support channel'' button to initiate Support
Mode.
\item After a brief period, the page will change to say the connection was
successful, if the connection was not successful, refer to the troubleshooting
section.
\item Do not close the web page, or click the ``Back'' or ``Refresh''
buttons. It is important to stay on this page until the Synercon Support
representative is finished providing support.
\item When the Synercon Support representative begins working on your FLA,
the screen on the device will change to indicate a member is preforming
work.
\item When the Synercon Support representative is finished, the unit will
display a message indicating the work is finished.
\item The FLA will automatically reboot.
\item After the device loads the Title Screen, then click the ``Close Support''
button to return the FLA to normal operational mode.
\end{enumerate}

\section{Installing PC Software}


\subsection{RP 1210 Drivers}

The FLA can opperate as an RP1210 device. This is used for connecting
OEM software to ECMs. A computer does not need the RP 1210 drivers
to access the FLA Preview, and the FLA Portal. You will need administrative
privileges to install these drivers.

To install the RP1210 Drivers, follow these steps
\begin{enumerate}
\item Connect to the FLA Preview discribed in the above section (Ether In-Field,
or In-Office methods will work)
\item Select the Download Drivers option from the menu at the top of the
web page
\item Select ``Download Drivers''
\item Once the file is downloaded, extract the .zip file to a folder
\item Open the folder, and double-click ``DPA4Plus\_RP1210.exe''
\item Follow the on-screen instructions to install the drivers
\end{enumerate}
After the drivers are installed, they can be tested with the following
steps
\begin{enumerate}
\item Make sure the FLA is powered on, and connected to an ECM (It is possible
to test the drivers without a live ECM, to do so see the special instructions)

\begin{enumerate}
\item Skip steps 2,3,and 4 if not connected to a live ECM
\item Begin a scan on the vehicle
\item After the scan has completed and you see the "Enable
RP1210 Mode" screen
\item Confirm the top line of the FLA read "CAN/J1939 Pass
Mode."
\item Search on your computer for the program "Adapter Validation
Tool" and run it
\item Select "MD/HD"
\item Set the Vendor to "DG1210xx - Dearborn Group RP1210"
(the xx will likely be 32)
\item Set the Device to "DG DPA 4 Plus USB"
\end{enumerate}

\begin{enumerate}
\item Skip to step 14 if the ECM does not support J1939
\end{enumerate}
\item Set Protocol to "J1939"
\item plug the USB cable into the FLA and the computer
\item Click "Run Test"
\item The text boxes labeled "RP1210 Status Window (PC-VDA)"
and "RP1210 Data Message Window (VDA-Vehicle)"
should turn green

\begin{enumerate}
\item If you are not connected to a live ECM, the "RP1210
Data Message Window" will turn red as there are no messages
\end{enumerate}
\item A text box will appear and will describe the results of the test.
\item Click "Quit" on the pop-up window titled
"Test Results Discussion"

\begin{enumerate}
\item If the ECM does not support J1708, the test is complete.If it does,
continue
\end{enumerate}
\item Change the protocol to J1708
\item On the FLA click the Green button to enable J1708 passthrough
\item Confirm the top line of the FLA read "J1708 Pass Mode
for"
\item Click "Run Test"
\item A text box will appear and will describe the results of the test
\item Click "Quit" on the pop-up window titled
"Test Results Discussion"
\item Close out of the adapter validation test
\end{enumerate}

\section{Preforming a download on a Heavy Vehicle}

This section provides a very high level overview of a typical download.
\begin{enumerate}
\item Plug in the FLA to the 9-Pin Diagnostics Connector
\item Turn on the ignition key to run
\item Run a vehicle Scan
\item Download Supported Data (if available)
\item Enable RP1210 Mode (Optional depending on the ECM and user's proceedure)

\begin{enumerate}
\item Connect the USB Cable
\item Run OEM Software
\end{enumerate}
\item Verify the data collected via the FLA Preview

\begin{enumerate}
\item If desired, repeat steps 3-6 until satisfied
\end{enumerate}
\item Shutdown the FLA
\item Reach a location that has a working Ethernet port
\item Connect the FLA to the Internet
\item Upload the data from the field
\item View the full report on the FLA Portal
\end{enumerate}

\section{Connecting to Engine Control Modules on a Bench}

When preforming bench downloads, it is recommended to use a Smart
Sensor Simulator to mimimize the chance of throwing additional faults
as compared to a traditional calibration harness. For more information
about Smart Sensor Simulators, visit \textbackslash{}url\{synercontechnologies.com\}.
Once connected, use the FLA in accordance to the recomended practices
outlined in the Field Guide and accompanying Quick Start guides.

The following are typical steps taken during a download.
\begin{enumerate}
\item Confirm the correct SSS is being used for the ECM
\item Connect the SSS to the ECM
\item Connect power to the SSS via the included power adapter
\item Press the Key Switch until the Key On lamp is iluminated, about two
seconds
\item Verify the Key On lamp stays iluminated for one second before proceeding.
If not, refer to the troubleshooting section
\item Plug in the FLA to the 9-Pin Diagnostics Connector located on the
SSS
\item Run a vehicle Scan
\item Download Supported Data (if available)
\item Enable RP1210 Mode (Optional depending on the ECM and user's proceedure)

\begin{enumerate}
\item Connect the USB Cable
\item Run OEM Software
\end{enumerate}
\item Verify the data collected via the FLA Preview

\begin{enumerate}
\item If desired, repeat steps 3-6 until satisfied
\item Skip to step 13 if there is a working Ethernet port where the bench
download is occuring
\end{enumerate}
\item Shutdown the FLA
\item Reach a location that has a working Ethernet port
\item Connect the FLA to the Internet
\item Upload the data from the field
\item View the full report on the FLA Portal
\end{enumerate}

\chapter{Troubleshooting}

\section{Vehicle Connection Problems}
\subsection{Do not get the Option to Scan Vehicle}
Symptopm:\\
\noindent\begin{minipage}{0.45\textwidth}% adapt widths of minipages to your needs
The FLA Title Screen does not allow the user to scan the vehicle
\end{minipage}%
\hfill%
\begin{minipage}{0.45\textwidth}
\includegraphics[width=\linewidth]{../media/fla_screens/title_screen_no_veh_net}
\end{minipage}\\
Solution:\\
\begin{itemize}
\item Check that the key is in the ON position, with the engine not running.
\item Check to see if the vehicle voltage is nominal, low voltage may prevent ECMs from responding.
\item Check all cables for secure connections.
\end{itemize}
If the above list does not solve the issue, the communication wires from the ECM to the diagnostic port may be damaged. If this is the case, the ECM modules may have to be removed and the operator may have to preform a bench download on the unit. The recomended practice is to use a Synercon Smart Sensor Simulator with the truck ECMs in order to help preserve the data on the modules.
\section{FLA Internet Problems}
\subsection{Cannot access the FLA Preview website}
Symptom:\\
The operator is not able to access the FLA Preview website.
Solution:\\
\begin{itemize}
\item navagate to the title screen and verify the FLA has an IP address. If the IP line reads 'cable unplugged', or 'finding IP address', see the FLA Does not have an IP Address section.
\item If the FLA is not connected directly to a computer (i.e. the FLA is plugged into an Ethernet port in an office, or home) ensure the DHCP service is off. If the IP address of the FLA is 10.0.0.1, the DHCP service is likely on. For more information about the DHCP service, see section~\ref{subsec:dhcp_services}.
\item Ensure the operator has typed in the IP address exactly as it appears on the FLA, do not add any extra characters, such as leading zeros before the numbers, or extra spaces.
\end{itemize}
\noindent\begin{minipage}{0.45\textwidth}% adapt widths of minipages to your needs
\begin{center}
\textbf{Correct:}\\[\baselineskip]
\end{center}
\includegraphics[width=\linewidth]{../media/fla_preview_screenshots/url_correct}
\begin{center}
\textbf{Incorrect:}\\[\baselineskip]
\end{center}
\includegraphics[width=\linewidth]{../media/fla_preview_screenshots/url_wrong_1}
\\[0.5\baselineskip]
\includegraphics[width=\linewidth]{../media/fla_preview_screenshots/url_wrong_2}
\\[0.5\baselineskip]
\includegraphics[width=\linewidth]{../media/fla_preview_screenshots/url_wrong_3}
\end{minipage}%
\hfill%
\begin{minipage}{0.45\textwidth}
\includegraphics[width=\linewidth]{../media/fla_screens/title_screen_no_veh_net}
\end{minipage}
\begin{itemize}
\item Ensure the FLA is on the same network as the computer used to access it. Depending on the configuration, if the coputer is connected to a wireless network, try to plug into an Ethernet cable near the port the FLA is plugged into.
\end{itemize}
If the above steps do not resolve the problem, the operator may need to contact the department IT services. If they are not able to resolve the issue, please contact Synercon support.

\subsection{FLA Does not Have an IP Address}
\subsubsection{IP says Ethernet Unplugged}
Symptom:\\
\noindent\begin{minipage}{0.45\textwidth}% adapt widths of minipages to your needs
The IP line of the FLA Title screen says Ethernet Unplugged
\end{minipage}%
\hfill%
\begin{minipage}{0.45\textwidth}
\includegraphics[width=\linewidth]{../media/fla_screens/title_screen_eth_unplug}
\end{minipage}\\
Solution:\\
\begin{itemize}
\item Check cables are plugged in all the way.
\item Reboot the FLA with the ethernet cable plugged in.
\item Verify the port is working with a computer.
\item Contact your deprartment IT services, and ensure the FLA is given the permissions to access the internet. Often this is done by 'whitelisting a mac address'. If the operator needs to provide the MAC of the FLA, this can be accessed via the configuation menu. For more information, see section~\ref{subsec:finding_mac}.
\end{itemize}
If the above steps do not solve the issue, the operator may need to contact Synercon support.
\subsubsection{IP says Finding IP Address}
Symptom:\\
\noindent\begin{minipage}{0.45\textwidth}% adapt widths of minipages to your needs
The IP line of the FLA Title screen says Ethernet Unplugged
\end{minipage}%
\hfill%
\begin{minipage}{0.45\textwidth}
\includegraphics[width=\linewidth]{../media/fla_screens/title_screen_eth_unplug}
\end{minipage}\\
Solution:\\
\begin{itemize}
\item If the FLA is being used in the field, and directly connected to a computer, enable the DHCP service. See section~\ref{subsec:dhcp_services}.
\item Reboot the FLA with the ethernet cable plugged in.
\item Verify the port is working with a computer.
\item Contact your deprartment IT services, and ensure the FLA is given the permissions to access the internet. Often this is done by 'whitelisting a mac address'. If the operator needs to provide the MAC of the FLA, this can be accessed via the configuation menu. For more information, see section~\ref{subsec:finding_mac}.
\end{itemize}
\section{FLA Boot Problems}
\subsection{FLA Does not Boot When Plugged in}
Symptom:\\
The FLA does nothing when plugged into a vehicle.
Solution:\\
\begin{itemize}
\item The power wires to the diagnostic port may have been damaged, use the cigarette lighter adapter in the FLA case to provide power to the FLA. It is recommended to cary a portable battery for powering the FLA via battery for such situations.
\end{itemize}
If the FLA will not boot with either the cigarette lighter adapter plugged into a know charged battery, or the wall adapter, please contact Synercon support.
\subsection{FLA Does not finish booting}
Symptom:\\
\noindent\begin{minipage}{0.45\textwidth}% adapt widths of minipages to your needs
The FLA reaches this screen, and does not proceed to the Title screen for more than 30 seconds.
\end{minipage}%
\hfill%
\begin{minipage}{0.45\textwidth}
\includegraphics[width=\linewidth]{../media/fla_screens/boot_wait}
\end{minipage}\\
Solution:\\
\begin{itemize}
\item Wait 30 seconds for the system to attempt to boot.
\item If the system does not boot, remove power sources, wait 30 seconds and then plug the FLA in. Note this is not recommended as a standard method to shut the FLA down, only do this when you have exhausted all other methods to shut the system down safely.
\end{itemize}
If the device continues to fail to boot for more than 2 times, please contact Synercon support.


\chapter{Definitions}

When dealing with heavy truck data, there are many acronyms, symbols
and phrases that are used to describe the data. This chapter presents
a list to help decipher some of the terms used for dealing with heavy
vehicle event data recorders.

The Forenisc Link Adapter (FLA) is a field rugged computer that is programmed to download data from heavy vehicles.
\begin{description}
\item [{ECMs:}] Electronic Control Modules. These are the ``computers''
that run heavy vehicles, they can contain data useful when investigating
a heavy vehicle.
\item [{OEM}] Original Engine Manufacturer.
\item [{The FLA}] the data extraction device device created by Synercon.
This is the blue box that the operator uses to comminicate with heavy
vehicles.
\item [{The FLA Portal}] the website fla.synercontechnologies.com . This
is where operators can view full reports
\item [{The FLA Preview}] the local website hosted by the FLA. This allows
operators to preview reports before they are uploaded to the FLA Portal.
\item [{Data Package}] all of the data extracted from a heavy vehicle,
as well as metadata about the extracted data, such as GPS coordinates,
and timestamps.
\item [{Preview Report}] a report that has not been processed by the FLA
Portal, these reports only have standards-based data, and some of
the non-standards-based data.
\item [{Full Report}] a report that is being viewed from the FLA Portal,
these reports will have all of the data the FLA system is able to
understand.
\item [{Standards-based-data}] Data that is defined by the SAE standards
J1587 and J1939.
\item [{Non-Standards-based-data}] Data that is not defined by the SAE
standards J1587 and J1939. i.e. Snapshot records for CAT ECMs. These
data elements are determined by the OEM of the ECMs
\item [{NIST}] National Institute for Standards and Technology (\href{http://www.nist.gov/}{http://www.nist.gov/})
\item [{RTC}] Real Time Clock
\end{description}

\chapter{Requirements}
\paragraph{  }
In order to take full advantage of the Synercon products, such as
the FLA, FLA Preview, and FLA Portal, the user will need to have access to additional software and ports such as those outlined in the following sections.


\section{Supported Web Browsers}
\paragraph{  }
In order for the FLA Preview and FLA Portal to work correctly, the
user will need to meet the following browser requirements.

\begin{description}
\item [{Internet Explorer}] at least version 10 (2012)
\item [{Firefox}] at least version 31 (2012)
\item [{Chrome}] at least version 31 (2013)
\item [{Safari}] at least version 7 (2013)
\item [{Opera}] at least version 29 (2015)
\end{description}

\section{Drivers}
\paragraph{  }
In order to use the FLA in RP 1210 mode, the operator will need to
install the drivers first. This will require administrative privilages.
The drivers are provided with the FLA.


\section{OEM Software}
\paragraph{  }
In order to extract data from ECMs that are not fully supported, the
operator will need the relative OEM software to aquire non-standards-based
data.


\chapter{Warnings}
\paragraph{  }
Synercon Technologies does not warrante any products against user
caused problems. In order to prevent any damage to the FLA and data
collected by the FLA, always follow these guidelines.


%\section{Powering the FLA}


\subsection*{Only use the included wall adapter or cigarette adapter to power
the FLA}

When not using a vehicle to power the FLA, only power the device using
the included power supplies. Never plug in a forign adapter, espeically
if the voltage exceeds 12 volts.


\subsection*{Never plug the FLA into a diagnostic port that has a vehicle voltage
above 26 volts}

When plugging in the FLA to a vehicle, ensure the voltage of the diagnostic
port does not exceed 26 volts. If in doubt, a multi-meter can be used
to verify voltage


\subsection*{Do not remove power from the FLA without first shutting down}

Unless specifically instructed to do so by a troubleshooting guide,
avoid removing power to the FLA without first shutting down. Use the
securing screws of the cables to secure all cables against accidental
dislodgement.


\subsection*{Do not expose the FLA to water}

The FLA is NOT waterproof, even powered off. Avoid dropping the unit
in puddles, or exposing the device to rainfall. Take care to ensure
a dry environment is avaliable before operating the FLA. If the unit
is exposed to water, even when powered off, the water can damage the
unit, and potentially any data stored on the unit. If exposed to water,
imediatly power the unit off and contact Synercon Support for instructions.


\subsection*{Do not expose to excess heat or cold}

Do not expose the FLA to temperatures exceeding 50C (122F), or below
0C (-32F).


\chapter{Synercon Training}

Synercon Technologies provides classes to train users in the proper
use of Synercon Products in investigating vehicles. For more information about classes, please contact XXX.

\begingroup
\renewcommand{\cleardoublepage}{}
\renewcommand{\clearpage}{}
\chapter{Contact Synercon Techonlogies}\label{chap:contact_syn}
\endgroup
\section*{Sales}
\paragraph{  }
If you need to extend an extending an FLA subscription, or are interested in puchasing additional FLAs, or SSSs, please contact our sales department at email\{sales@synercontechnologies.com\}
\section*{Support}
\paragraph{  }
While we strive to provide the hightest quality products, if you have
an issue with our products and need to contact us, please send an
email to \textbackslash{}email\{support@synercontechnologies.com\}.

When requesting support, please include the following:

1) the product (i.e. FLA, or SSS)

2) the serial number (it is printed on the bottom of all Synercon
Products)

3) A detailed explanation of the issue

4) Any steps you have already taken to solve the issue

5) {[}optional{]} any screen shots, or pictures if you believe they
will help us understand the problem.

We will attempt to respond to you as quickly as possible.

\printindex{}


%\section{Specifications}
%
%The Forensic Link Adapter is based on a powerful processor based on
%the Texas Instruments AM335x 1GHz ARM\textregistered{} Cortex-A8 Processor.
%It has 512MB of DDR3 RAM, 16 GB of internal storage, the NEON floating-point
%accelerator and 2 PRU 32-bit microcontrollers. The device has 2 Controller
%Area Network (CAN) channels, 2 J1708 Serial Channels, 100/10 Mbps
%Ethernet, USB 2.0, and a full implementation of the DG tech DPA4+
%for RP1210 Compliance. The device has a GPS receiver, Real Time Clock,
%an accelerometer, and a rate Gyro. The display uses Organic LEDs for
%better sunlight readability. It runs embedded Linux and serves its
%own website.

\end{document}
